% Generated draft: requires mathematical and source review.
\documentclass[11pt]{article}
\usepackage[margin=25mm]{geometry}
\usepackage{fontspec}
\setmainfont{Latin Modern Roman}
\usepackage{amsmath,amssymb,mathtools,unicode-math}
\setmathfont{Latin Modern Math}
\usepackage{enumitem,needspace,xurl,hyperref}
\setlength{\parindent}{0pt}
\setlength{\parskip}{5pt}
\setlength{\emergencystretch}{3em}
\newcommand{\N}{\mathbb{N}}
\newcommand{\Z}{\mathbb{Z}}
\newcommand{\Q}{\mathbb{Q}}
\newcommand{\R}{\mathbb{R}}
\newcommand{\C}{\mathbb{C}}
\newcommand{\F}{\mathbb{F}}
\newcommand{\A}{\mathbb{A}}
\newcommand{\PP}{\mathbb{P}}
\newcommand{\E}{\mathbb{E}}
\newcommand{\eps}{\varepsilon}
\newcommand{\dd}{\,\mathrm{d}}
\begin{document}
\section*{UnsolvedMath ID: 20000444}
\textbf{Source problem code:} catalogue entry
\subsection{Definitions and mathematical statement}
This record is a historical status remark on Question 11, not the statement of a new unresolved problem. For a field \(k\), write \(\operatorname{cd}(k)\) for its Galois cohomological dimension, subject to the convention used in the cited sources. For a nonempty proper \(k\)-variety \(X\), its index is
\[
\operatorname{ind}(X)=\gcd\{[k(P):k]:P\text{ a closed point of }X\}
=\gcd\{[L:k]:L/k\text{ finite},\ X(L)\ne\varnothing\}.
\]
A zero-cycle is an integral linear combination of closed points, with degree given by the corresponding combination of residue-field degrees. Thus index one means existence of a degree-one zero-cycle; its coefficients need not be nonnegative and it need not be a rational point. Rational connectedness refers to the geometric property that two general points over an algebraic closure can be joined by a rational curve.

The supplied source reports the following two counterexamples, over potentially different fields.
\begin{enumerate}
\item J. Ax constructed a field \(k\) with \(\operatorname{cd}(k)=1\) and a smooth degree-five hypersurface \(X\subset\PP^9_k\) such that
\[
X(k)=\varnothing,\qquad \operatorname{ind}(X)=1.
\]
The source notes that rational connectedness of such hypersurfaces was established later. Thus the rational-point assertion attached to Question 11 was already false, even for a rationally connected variety of index one.
\item J.-L. Colliot-Th\'el\`ene and D. Madore constructed a field \(k'\) with \(\operatorname{cd}(k')=1\) and a smooth cubic surface \(Y\subset\PP^3_{k'}\) satisfying
\[
\operatorname{ind}(Y)=3.
\]
In particular it has no degree-one zero-cycle. The source states that this gives a negative answer to a question of Kato and Kuzumaki (1986); the attached aid identifies the condition as \(C_1^0\).
\end{enumerate}
The precise original wording of Question 11, the fields' characteristics, and the convention for cohomological dimension are not reproduced in this continuation. They should not be filled in by conjecture.

The attached research aid records an auxiliary comparison: for a nonempty proper variety \(Z/k\) and every finite extension \(L/k\), it claims
\[
\operatorname{ind}(Z_L)\mid\operatorname{ind}(Z),\qquad
\operatorname{ind}(Z)\mid[L:k]\operatorname{ind}(Z_L).
\]
Consequently, if \(\operatorname{ind}(Z)=p\) is prime and \(p\nmid[L:k]\), its index remains \(p\) after this extension. The aid also uses that a degree-\(p\) hypersurface has index dividing \(p\), and hence index \(1\) or \(p\). At Milnor degree zero, \(K^M_0(k)=\Z\) and the norm from a finite extension is multiplication by its degree. Accordingly, the subgroup generated by closed-point norms is
\[
N_0(Z/k)=\operatorname{ind}(Z)\Z.
\]
This distinguishes the Ax example, whose obstruction is not the absence of a degree-one zero-cycle, from the cubic's obstruction \(N_0(Y/k')=3\Z\). These are supplied research-aid comparisons, not a new open target. The aid additionally mentions an explicit Ax zero-cycle \([P_3]-[P_2]\); no definition or construction of the alleged degree-three and degree-two points is supplied, so that particular representation cannot be inferred merely from index one.
\subsection{Short English statement}
This is a status note reporting two established counterexamples: an index-one rationally connected quintic without a rational point, and an index-three cubic surface, each over a field of cohomological dimension one. The second has no degree-one zero-cycle. The registry's open label does not turn this continuation into a fresh open question.
\subsection{Sources}
\begin{itemize}
\item AIM, \emph{Rational and integral points on higher dimensional varieties}, remark (iv) on printed p. 50, immediately before Question 12, according to the supplied placement verification: \url{https://aimath.org/WWN/qptsurface2/qptsurface2.pdf}.
\item J. Ax, \emph{Bulletin of the American Mathematical Society} \textbf{71} (1965), p. 717, as cited in the source; no article title is supplied.
\item J.-L. Colliot-Th\'el\`ene and D. Madore, \emph{Journal of the Institute of Mathematics of Jussieu} \textbf{3} (2004), pp. 1--16, as cited in the source; no article title is supplied. The source mentions the Kato--Kuzumaki question (1986) without a full citation.
\item Local registry \texttt{data/unsolvedmath/problems.json}, permanent ID \(20000444\), exact source code \texttt{catalogue entry}; canonical \texttt{aim-algebraic-number-theory-notes.json}, notes[95]. The attached research attempt is explicitly machine-generated, not peer reviewed, and requires independent verification before citation.
\end{itemize}
\end{document}
